% !TeX root = ../../Book1.tex
% 纲要标题：### 4.4 群作用、置换嵌入与子群正规性
% 纲要位置：计划纲要.md，4.4 节。

\section{群作用、置换嵌入与子群正规性}
\label{b1:sec:0404}

一个群的作用除了帮助计数，也能把抽象群放进熟悉的置换群中。作用忠实时得到置换嵌入；作用非平凡而核也非平凡时，核给出非平凡真正规子群。左乘作用保证嵌入总能实现，陪集作用常把置换次数降到较小的指数。下面几个例子说明，选择作用集合本身就是结构分析的一步。

% 本节正文按下列原纲要要点撰写。

% 纲要内容（仅作写作计划，尚非教材正文）：
%
% * Cayley 定理与正则作用
% * 低指数子群与正规性
% * 通过不同作用比较群的结构
%
% ---
%

\subsection{Cayley 定理与正则作用}
\label{b1:subsec:0404:01}

在抽象群里，乘法表看似没有几何背景；但每个元素都给出对整个群的左乘置换。这一简单观察产生\term[cayley-dingli]{Cayley 定理}[Cayley theorem]。\footnote{凯莱（Arthur Cayley，1821-1895），英国数学家，发展矩阵代数，也研究群论与不变量理论。}

\begin{theorem}[Cayley]\label{b1:thm:cayley}
任意群 \(G\) 同构于某个集合上的置换群。若 \(|G|=n<\infty\)，则 \(G\) 同构于 \(S_n\) 的一个子群。
\end{theorem}
\begin{proof}
令 \(G\) 在其底集合上左乘，得到同态
\[
\lambda:G\longrightarrow S_G,\qquad
\lambda(g)(x)=gx.
\]
若 \(\lambda(g)=\lambda(h)\)，则在 \(x=1\) 处取值得到 \(g=h\)。因此 \(\lambda\) 单射，左乘作用忠实。把目标限制到像 \(\lambda(G)\)，这个同态便成为双射，所以 \(G\cong\lambda(G)\leq S_G\)。有限情形给 \(G\) 的 \(n\) 个元素编号，就把 \(S_G\) 识别为 \(S_n\)。
\end{proof}

\begin{definition}\label{b1:def:regular-action}
设群 \(G\) 作用在非空集合 \(X\) 上。如果对任意 \(x,y\in X\)，都恰有一个 \(g\in G\) 满足 \(gx=y\)，则将此作用称为\term[zhengze-zuoyong]{正则作用}[regular action]。
\end{definition}

正则作用恰好是自由而传递的作用。若作用正则，则对任意 \(x,y\in X\) 都存在 \(g\) 使 \(gx=y\)，这就是传递性；取 \(y=x\) 时，单位元已经固定 \(x\)，唯一性便要求没有其他元素固定它，所以作用自由。

\ProseParagraphBreak
反过来，设作用自由而传递。对任意 \(x,y\in X\)，传递性先保证至少有一个 \(g\) 使 \(gx=y\)。若另有 \(h\) 也满足 \(hx=y\)，则 \(h^{-1}gx=x\)，自由性给出 \(h^{-1}g=1\)，即 \(g=h\)。因此传递性负责存在，自由性负责唯一，合在一起就是正则性。

\ProseParagraphBreak
设 \(G\) 为群。群 \(G\) 在其底集合上由 \(g\cdot x=gx\) 给出的作用称为\term[zuo-zhengze-zuoyong]{左正则作用}[left regular action]；方程 \(gx=y\) 的唯一解为 \(g=yx^{-1}\)。例如，令 \(C_3=\{1,a,a^2\}\)，并把这三个元素依次编号为 \(1,2,3\)。左乘给出
\[
\lambda(1)=1,\qquad
\lambda(a)=(1\,2\,3),\qquad
\lambda(a^2)=(1\,3\,2).
\]
这里 \(\lambda(a)^2=\lambda(a^2)\)、\(\lambda(a)^3=1\)，群乘法在这个具体的置换群中表现为置换复合。

\ProseParagraphBreak
任一正则作用选定一点 \(x\) 后，其稳定子群为 \(G_x=\{1\}\)。\cref{b1:prop:transitive-coset}给出它与 \(G/\{1\}\) 的等变同构，也就是与 \(G\) 的左乘作用的等变同构。正则作用所需点数可能偏大：\(S_3\) 的正则作用在六个点上，而它的自然忠实作用只用三个点。寻找较小的忠实作用，是在保留结构信息的同时减少置换计算量的一种方法。

\subsection{低指数子群与正规性}
\label{b1:subsec:0404:02}

子群 \(H\leq G\) 给出左陪集作用，从而给出同态
\(\rho_H:G\rightarrow S_{G/H}\)，其中 \(\rho_H(a)(gH)=agH\)。一个元素固定全部陪集，意味着它同时满足每个陪集的稳定条件。

\ProseParagraphBreak
从群结构来看，我们希望在 \(H\) 中找到一个对整个 \(G\) 正规、而且尽可能大的子群。若 \(N\trianglelefteq G\) 且 \(N\leq H\)，则对每个 \(g\in G\) 都有
\[
N=gNg^{-1}\leq gHg^{-1}.
\]
因此这样的 \(N\) 必须同时包含在 \(H\) 的所有共轭子群中。这解释了为什么要取所有共轭子群的交；下面证明这个交恰是陪集作用的核。

\begin{proposition}[陪集作用的正规核]\label{b1:prop:coset-action-core}
设 \(G\) 为群，\(H\leq G\)，并令 \(\rho_H:G\rightarrow S_{G/H}\) 为 \(G\) 在左陪集集合 \(G/H\) 上由左乘诱导的作用同态。则
\begin{equation}\label{b1:eq:coset-action-core}
\ker\rho_H=\bigcap_{g\in G}gHg^{-1}
=\operatorname{Core}_G(H).
\end{equation}
这个子群称为 \(H\) 在 \(G\) 中的\term[zhenggui-he]{正规核}[core]。它是所有满足 \(N\trianglelefteq G\) 且 \(N\leq H\) 的子群 \(N\) 中最大的一个。陪集作用忠实，当且仅当 \(\operatorname{Core}_G(H)=\{1\}\)。
\symidx{\(\operatorname{Core}_G(H)\)：子群在给定群中的正规核}
\end{proposition}
\begin{proof}
元素 \(a\) 固定陪集 \(gH\)，当且仅当 \(agH=gH\)，即 \(g^{-1}ag\in H\)，也就是 \(a\in gHg^{-1}\)。对所有陪集同时要求这个条件，就得到\eqref{b1:eq:coset-action-core}。

正规核作为同态的核正规于 \(G\)，并且在交中取 \(g=1\) 可见它包含于 \(H\)。若 \(N\trianglelefteq G\) 且 \(N\leq H\)，则对每个 \(g\in G\)，有
\(N=gNg^{-1}\leq gHg^{-1}\)，故 \(N\leq\operatorname{Core}_G(H)\)。这就证明了最大性，其中的正规性始终相对于 \(G\)。最后，\cref{b1:def:faithful-action}说明作用忠实恰好等价于其核平凡。
\end{proof}

例如在 \(G=S_3\) 中取 \(H=\bigl\langle(1\,2)\bigr\rangle\)。用 \((2\,3)\) 共轭得到 \(\bigl\langle(1\,3)\bigr\rangle\)，这两个二阶子群的交为 \(\{1\}\)，所以 \(\operatorname{Core}_{S_3}(H)=\{1\}\)。这里 \(H\) 自身正规于 \(H\)，却不包含任何非平凡的 \(S_3\)-正规子群；两种正规性不能混用。

\begin{corollary}[子群共同固定的陪集]\label{b1:cor:fixed-cosets-conjugate-containment}
设 \(G\) 为群，\(P,H\leq G\)。将 \(G\) 在左陪集集合 \(G/H\) 上的左乘作用限制到 \(P\)，并以 \((G/H)^P\) 表示被 \(P\) 中每个元素固定的陪集所成集合。则
\[
(G/H)^P
=\{\,gH\in G/H\mid g^{-1}Pg\leq H\,\}
=\{\,gH\in G/H\mid P\leq gHg^{-1}\,\}.
\]
\end{corollary}
\begin{proof}
左陪集作用中，基点 \(H\) 的稳定子群是 \(H\)；把基点换到 \(gH\)，同一轨道中的稳定子群共轭关系给出 \(G_{gH}=gHg^{-1}\)。因此 \(P\) 共同固定 \(gH\)，就是要求 \(P\) 的每个元素都属于这个点的稳定子群，即 \(P\leq G_{gH}=gHg^{-1}\)。用代表元逐项核验，得到
\[
\begin{aligned}
gH\in(G/H)^P
&\iff pgH=gH\quad\text{对所有 }p\in P\\
&\iff g^{-1}pg\in H\quad\text{对所有 }p\in P\\
&\iff g^{-1}Pg\leq H\\
&\iff P\leq gHg^{-1}.
\end{aligned}
\]
这组等价关系同时给出两个集合等式。
\end{proof}

在\cref{b1:prop:coset-action-core}中，先固定群元素，再要求它保持全部陪集，所得条件描述作用核；在上述推论中，先固定一个陪集，再要求子群的全部元素都保持它，所得条件描述共轭包含。两处使用的是同一个等式 \(pgH=gH\)，量词方向却不同。

\ProseParagraphBreak
仍取 \(G=S_3\)、\(H=\bigl\langle(1\,2)\bigr\rangle\)，并令 \(P=\bigl\langle(2\,3)\bigr\rangle\)。三个左陪集为 \(H,(1\,3)H,(2\,3)H\)。\(P\) 交换第一个与第三个陪集，却固定 \((1\,3)H\)，因为
\[
(1\,3)^{-1}P(1\,3)=\bigl\langle(1\,2)\bigr\rangle=H.
\]
因此 \((G/H)^P=\{(1\,3)H\}\)。这个二阶群作用在三个陪集上，\(\lvert(G/H)^P\rvert=1\equiv3=\lvert G/H\rvert\pmod2\)，也直接体现了\cref{b1:lem:p-action-fixed}的不动点同余。

\ProseParagraphBreak
当有限群 \(G\) 的子群 \(H\) 有素数指数 \(p\) 时，陪集作用给出 \(G\rightarrow S_p\)。要证明 \(H\) 正规，可以尝试把它识别为这个作用的核 \(K\)。我们已知 \(K\leq H\)；若能证明作用的像恰有 \(p\) 个元素，第一同构定理就给出 \([G:K]=p=[G:H]\)，进而 \([H:K]=1\)，所以 \(H=K\)。下面的最小素因子条件将限制像的阶，使这一步成立。

\begin{theorem}\label{b1:thm:least-prime-index-normal}
设 \(G\) 为有限群，\(p\) 是 \(|G|\) 的最小素因子。若 \(H\leq G\) 且 \([G:H]=p\)，则 \(H\trianglelefteq G\)。
\end{theorem}
\begin{proof}
令 \(\rho_H:G\rightarrow S_{G/H}\cong S_p\) 为左陪集作用同态，并令 \(K=\ker\rho_H\)。由\eqref{b1:eq:coset-action-core}，\(K\leq H\)，且核 \(K\) 正规。\cref{b1:thm:lagrange}分别用于 \(K\leq H\) 与 \(H\leq G\)，得到
\[
|G/K|=\frac{|G|}{|K|}
=\frac{|G|}{|H|}\frac{|H|}{|K|}=p[H:K].
\]
因此 \(p\bigm\vert|G/K|\)。接下来分别约束这个阶的素因子及其出现次数。

首先，\(|G/K|=|G|/|K|\) 整除 \(|G|\)，所以它的每个素因子都不小于 \(p\)，这里使用了 \(p\) 是 \(|G|\) 的最小素因子。

其次，\cref{b1:thm:first-isomorphism}给出 \(G/K\cong\operatorname{im}\rho_H\leq S_p\)。对这个有限子群使用\cref{b1:thm:lagrange}，其阶整除 \(p!\)，所以每个素因子都不大于 \(p\)。结合前一步，\(|G/K|\) 的素因子只能是 \(p\)，故 \(|G/K|=p^a\)，其中 \(a\geq1\)。

最后，\(1,2,\ldots,p\) 中只有 \(p\) 被 \(p\) 整除，且只含一个因子 \(p\)，所以 \(p!\) 中素因子 \(p\) 恰出现一次。由 \(p^a\mid p!\)，得到 \(a=1\)，即 \(|G/K|=p\)。代回前面的指数式，\([H:K]=1\)，故 \(H=K\) 正规。
\end{proof}

\cref{b1:prop:normal-basic-examples}已经证明，任意群中的指数为 \(2\) 的子群都正规。该处的直接证明只用到左、右陪集各有两个，因而不要求群有限。指数为其他素数的子群则未必正规，\(S_3\) 的换位子群有指数 \(3\)，但不是正规子群；这里 \(3\) 不是群阶的最小素因子。

\subsection{群作用与结构比较}
\label{b1:subsec:0404:03}

选择作用时，先决定希望看见什么。下面比较同一个群 \(G\) 的四种作用，其中 \(\mathcal S(G)\) 表示全部子群组成的集合，\(H\leq G\) 是指定子群。
\begin{center}
\begin{tabular}{llll}
\toprule
作用与集合 & 轨道 & 点稳定子群 & 作用核 \\
\midrule
\(G\)，左乘 & 整个 \(G\) & \(\{1\}\) & \(\{1\}\) \\
\(G\)，共轭 & 元素共轭类 & \(C_G(x)\) & \(Z(G)\) \\
\(\mathcal S(G)\)，共轭 & 共轭子群族 & \(N_G(L)\) & \(\bigcap_{L\leq G}N_G(L)\) \\
\(G/H\)，左乘 & 整个 \(G/H\) & \(H\) & \(\operatorname{Core}_G(H)\) \\
\bottomrule
\end{tabular}
\end{center}
第二行在元素 \(x\) 处取稳定子群，第三行在子群 \(L\) 处取稳定子群，最后一行则在基点 \(H\) 处取稳定子群；若换到 \(gH\)，它变为 \(gHg^{-1}\)。对子群集合的共轭作用之核未必等于中心：\(Q_8\) 的每个子群都正规，所以该核为整个 \(Q_8\)，而对元素集合的共轭作用之核为 \(Z(Q_8)=\{1,-1\}\)。陪集作用的核是指定子群的正规核，能把小指数条件转成低次置换群的限制。这些作用提供的轨道、稳定子群和核不同，应随所要解决的问题选择。

\ProseParagraphBreak
以 \(S_4\) 为例，我们想找一个非平凡真正规子群。若让它在一个较小集合上非平凡作用，作用核便自动正规；若像的阶又小于 \(24\)，核就不可能平凡。四个点恰有三种分成两个无序对的方式，这给出了一个三点作用的候选：
\[
\omega_1=12\mid34,\qquad
\omega_2=13\mid24,\qquad
\omega_3=14\mid23.
\]
令 \(\Omega=\{\omega_1,\omega_2,\omega_3\}\)。群 \(S_4\) 按逐个置换字母的方式作用在 \(\Omega\) 上，例如
\[
\sigma\cdot(12\mid34)
=\sigma(1)\sigma(2)\mid\sigma(3)\sigma(4).
\]
竖线两侧各表示一个无序二元子集，这两个子集之间也不区分次序。置换保留这种划分，复合后各字母的像与逐次作用相同，所以得到同态 \(\rho:S_4\rightarrow S_\Omega\cong S_3\)。

\ProseParagraphBreak
换位 \((12)\) 固定 \(\omega_1\)、交换 \(\omega_2,\omega_3\)；换位 \((13)\) 固定 \(\omega_2\)、交换另外两点。因此
\[
\rho((12))=(\omega_2\,\omega_3),\qquad
\rho((13))=(\omega_1\,\omega_3).
\]
用第一个换位共轭第二个，得到第三个换位 \((\omega_1\,\omega_2)\)。三个换位生成三点集上的全部置换，所以 \(\operatorname{im}\rho=S_3\)。这里所选的两个换位不生成 \(S_4\)，但它们的像已经生成整个目标群，这就足以证明 \(\rho\) 满射。

\ProseParagraphBreak
\cref{b1:thm:first-isomorphism}给出 \(S_4/\ker\rho\cong S_3\)，因此 \(|\ker\rho|=|S_4|/|S_3|=24/6=4\)。逐个代入可见
\[
V_4=\bigl\{1,(12)(34),(13)(24),(14)(23)\bigr\}
\]
中的四个元素固定全部三种划分，故 \(\ker\rho=V_4\)。这个群称为\term[klein-siyuan-qun]{Klein 四元群}[Klein four-group]。\footnote{克莱因（Felix Klein，1849-1925），德国数学家，提出以变换群研究几何的埃尔朗根纲领，并研究自同构与群作用。}它的三个非单位元素都为二阶。取其中两个不同元素 \(a,b\)，消去律表明 \(ab\) 不同于 \(1,a,b\)，故四个元素恰为 \(1,a,b,ab\)。又因 \(a,b,ab\) 都为二阶，
\[
ab=(ab)^{-1}=b^{-1}a^{-1}=ba.
\]
将 \(C_2\times C_2=(\mathbb Z/2\mathbb Z)^2\) 配备逐坐标加法，定义
\[
(\mathbb Z/2\mathbb Z)^2\longrightarrow V_4,\qquad
(\bar i,\bar j)\longmapsto a^ib^j.
\]
这里 \(\bar i,\bar j\) 分别表示整数 \(i,j\) 的模 \(2\) 类。关系 \(a^2=b^2=1\) 保证映射与代表整数的选择无关，\(ab=ba\) 保证它保持运算；四个像两两不同，故它是双射，给出 \(V_4\cong C_2\times C_2\)。
\symidx{\(V_4\)：Klein 四元群}

\ProseParagraphBreak
于是 \(V_4\trianglelefteq S_4\) 且 \(S_4/V_4\cong S_3\)。上述三种划分也就是正四面体的三对相对边；现在还可以从这个作用中确定稳定子群及其正规核。

\ProseParagraphBreak
取 \(H=(S_4)_{\omega_2}\)。由于像为整个 \(S_3\)，\(\Omega\) 是一个传递 \(S_4\)-集合。\cref{b1:thm:orbit-stabilizer}给出
\[
[S_4:H]=3,\qquad |H|=\frac{24}{3}=8.
\]
按正方形的顶点编号，\(r=(1\,2\,3\,4)\) 交换 \(\{1,3\}\) 与 \(\{2,4\}\)，\(s=(2\,4)\) 则分别保持这两个子集，所以 \(r,s\) 都固定 \(\omega_2\)。
由\cref{b1:prop:dihedral-normal-form}及正方形编号的说明，\(\langle r,s\rangle\) 正是阶为 \(8\) 的二面体群。因此 \(\langle r,s\rangle\leq H\) 且两者同阶，得到
\[
H=\bigl\langle(1\,2\,3\,4),(2\,4)\bigr\rangle\cong D_8.
\]
\cref{b1:prop:transitive-coset}给出等变双射
\[
S_4/H\longrightarrow\Omega,\qquad gH\longmapsto g\omega_2.
\]
一个元素固定全部陪集，当且仅当它固定对应的全部划分，所以这个陪集作用的核仍为 \(V_4\)。由\cref{b1:prop:coset-action-core}，
\[
\operatorname{Core}_{S_4}(H)=V_4.
\]
八阶子群 \(H\) 保持一个指定划分，四阶正规核 \(V_4\) 则保持全部三个划分；\(S_4/H\) 在这里是三元素陪集集合，\(S_4/V_4\) 才是与 \(S_3\) 同构的商群。

\ProseParagraphBreak
有时单个作用不忠实，但几个作用放在一起可以忠实。设 \(m\geq1\) 为整数，\(G\) 分别作用在 \(X_1,\ldots,X_m\) 上，令
\[
X=\coprod_{i=1}^{m}X_i
=\{\,(i,x)\mid 1\leq i\leq m,\ x\in X_i\,\},
\qquad g\cdot(i,x)=(i,gx).
\]
这称为不交并上的作用。它的核是各作用核的交，因为固定 \(X\) 的每个点，恰好等于固定每个 \(X_i\) 的全部点。

\ProseParagraphBreak
与此相比，笛卡尔积上的作用为
\[
g\cdot(x_1,\ldots,x_m)=(gx_1,\ldots,gx_m).
\]
作用公理在每个坐标上分别成立。不交并作用保留分块标记 \(i\)，不会把一个 \(X_i\) 中的点送到另一块；积作用则让同一个 \(g\) 同时作用于全部坐标。当各 \(X_i\) 都非空时，积作用的核也等于各原作用核的交，因为固定全部元组就要求它固定每个坐标集合的全部点。\footnote{若某个 \(X_i\) 为空，整个笛卡尔积为空，其作用核为 \(G\)，不一定等于各原作用核的交。}两个构造的核可以相同，轨道结构却可以不同。

\ProseParagraphBreak
例如取 \(G=(\mathbb Z/2\mathbb Z)^2\)，让第一坐标交换 \(\{1,2\}\)，第二坐标交换 \(\{3,4\}\)。对 \(a,b\in\{0,1\}\)，用 \(\bar a,\bar b\) 表示其模 \(2\) 余类，则四点不交并作用的同态 \(\rho:G\rightarrow S_{\{1,2,3,4\}}\) 为
\[
\rho(\bar a,\bar b)=(1\,2)^a(3\,4)^b.
\]
两个二点作用的核分别由第二坐标、第一坐标组成；它们的交为单位，故 \(\rho\) 忠实。这个作用有两个轨道。若把相同的两个二点作用用于笛卡尔积，则得到四点上的正则作用：每对坐标都可以由唯一的 \((\bar a,\bar b)\) 送到任意目标坐标对。因此积作用只有一个轨道；两种构造的核相同，轨道却不同。\cref{b1:ex:04-8}进一步计算积作用的点稳定子群，并给出另一个轨道分解的例子。

\begin{remark}\label{b1:rem:action-structural-use}
设非平凡群 \(G\) 在一个 \(m\) 元集合上作用，记作用同态为 \(\rho:G\rightarrow S_m\)，并令 \(K=\ker\rho\)。由\cref{b1:thm:first-isomorphism}，
\[
G/K\xrightarrow{\sim}\rho(G)\leq S_m.
\]
当前作用的核有三种可能：
\[
\begin{array}{c|l}
K=\{1\} & \rho\text{ 给出 }G\text{ 的置换嵌入}\\
\{1\}<K<G & K\text{ 是非平凡真正规子群}\\
K=G & \rho\text{ 是平凡作用}
\end{array}
\]
因此，\(\{1\}<K<G\) 当且仅当作用非忠实且非平凡。若只知 \(K\ne\{1\}\)，还须检查 \(K\ne G\)：例如素数阶群 \(C_p\) 在单点集上的作用之核为整个 \(C_p\)，由此并未得到非平凡真正规子群。要由当前作用得到 \(G\) 到 \(S_m\) 的嵌入，必须证明 \(K=\{1\}\)；当作用非平凡时，证明其核非平凡便得到非平凡真正规子群。下一章研究 Sylow 子群上的共轭作用时，也将分别检查核与像。
\end{remark}
